\section{The Chatnotes groupoid and motivic interval: exact maps and source boundaries}
\label{sec:chatnotes-groupoid-motivic}

This section audits physical lines 11088--14011 of the directly read
17,437-line \emph{$Z_n$ Symmetries} export.  It closes the only interval of
that file which lacked a direct-locator audit.  The chronology consists of
four successive expansions:
\begin{enumerate}[label=\textup{(\roman*)}]
\item lines 11088--11720 attach the $1$-motive and Nori sources and propose a
  symbolic/groupoid ladder and a period ladder;
\item lines 11721--12206 repeat the grading and period formulas and replace
  the unnamed groupoid by an inverse-edge shift groupoid;
\item lines 12207--12660 attach a much broader operator and motivic
  vocabulary without constructing the proposed comparison maps; and
\item lines 12677--14011 discover genuine partial homeomorphisms of
  $\mathbb Z_3$, then incorrectly identify their compact operator model with
  the positive-integer path algebra and promote that identification to a
  larger arithmetic and equilibrium package.
\end{enumerate}
The chronology contains real mathematics.  Its two torus coordinate systems,
its $3$-adic branch maps, and its logarithmic and arctangent formulas can all
be made exact.  The exact results require maps which the raw response did not
supply, and those independently proved repairs are explicitly separated from
the source assertions below.

The controlling raw file has 538,005 bytes and SHA-256
\nolinkurl{9a80f53764a9da42c0f8e0b6d19ce61f978980a448a683b9c45f04038716160d}.
The locator report is
\path{intake/chatnotes/ZN-SYMMETRIES-GROUPOID-MOTIVIC-DIRECT-LOCATOR-AUDIT-06.md}.

\subsection{What is already proved before this interval}

The raw grading
\[
 \deg(p)=(\ell(p),A(p)-\ell(p))
\]
and its torus action at lines 11130--11265 are correct in the target-left
category algebra already defined in
Section~\ref{sec:chatnotes-weighted-path-algebra}.  The finite maps at lines
11269--11298 repeat the maps audited in
Theorem~\ref{thm:finite-residue-semigroup-normal-form} and
\ref{prop:weighted-typed-representations}; the raw text again erases object
idempotents and again calls the nonempty-word image a monoid before adjoining
an identity.  No second finite representation is admitted from that
repetition.

The full countable-word pressure calculation at lines 11302--11354 and
13637--13684 is the calculation already proved in
Section~\ref{sec:symbolic-completion-geometric-boundaries}.  It applies to
the full label shift.  The raw response does not verify the hypotheses needed
to transfer a Gibbs or equilibrium theorem to the positive-integer image.

The groupoid arrow at lines 11358--11391,
\[
 R[\mathcal E_C]\leadsto\mathcal G_C\leadsto C^*(\mathcal G_C),
\]
does not declare a topology, Haar system, full or reduced completion, or an
algebra map.  The inverse-edge version at lines 11813--11885 declares a path
space and shift but still gives no topology and no proof that the shift is a
local homeomorphism.  The later $3$-adic formulas supply a different exact
route; it must not be identified with either of those arrows without the maps
proved below.

\subsection{The literature actually used}

Deaconu begins with the set of triples attached to a map and proves his
topological theorem when the map is a covering of a compact Hausdorff space
\cite[printed pp.~1779--1781]{Deaconu1995}.  The positive-odd set with the
discrete topology is not compact and $C$ is not a self-covering, so that
theorem is not invoked.  Its set-theoretic groupoid formula is used as a
reference point; the required discrete topology and Haar assertions are
proved directly here.

Huber--W\"ustholz define a $1$-motive $[L\to G]$ and its morphisms as maps of
complexes, compute $\log\alpha$ from the Kummer motive
$[\mathbb Z\xrightarrow{1\mapsto\alpha}\mathbf G_m]$, and identify all
$1$-periods with periods of $1$-motives
\cite[Definition~8.1, Section~10.2, and Proposition~12.10]{HuberWustholz2022}.
Their Theorem~10.4 explicitly uses characters
$(x,y)\mapsto x^ny^m$ of $\mathbf G_m^2$.
Huber--M\"uller-Stach define a Nori vertex as a triple $(X,Y,i)$ and permit
exactly two primitive edge types: a contravariant edge from an actual map of
pairs and a coboundary edge from an actual nested triple
\cite[Definition~9.1.1]{HuberMullerStach2025}.  Concatenation of integration
paths is not a third edge type.

Barbieri-Viale--Kahn define a Deligne $1$-motive and its morphisms by a
commutative square in source TeX lines 728--750, prove the fully faithful
curve-generated embedding in source lines 1736--1751, and define $LAlb$ in
source lines 3796--3800 \cite{BarbieriVialeKahn2016}.  Their source provides
the general machinery, not the particular point $(2,3)$, the character
$x^Ay^{-\ell}$, or a curve realizing that object.  The concrete squares below
are therefore independent constructions.

\subsection{The exact positive-integer groupoid}

For $n\geq0$ write
\[
 S_na(x)=\sum_{j=0}^{n-1}a(C^jx),\qquad S_0a(x)=0.
\]
Define
\begin{equation}
\label{eq:critical-arithmetic-groupoid}
 \mathcal G_{\mathrm{arith}}
 =\{(x,k,y):x,y\in\mathcal O,\ \exists m,n\geq0,\
 k=m-n,\ C^m(x)=C^n(y)\}.
\end{equation}
Put $r(x,k,y)=x$, $s(x,k,y)=y$,
\[
 (x,k,y)(y,l,z)=(x,k+l,z),
 \qquad (x,k,y)^{-1}=(y,-k,x),
\]
and give the set the discrete topology.

\begin{theorem}[independent arithmetic groupoid repair]
\label{thm:critical-arithmetic-groupoid}
The object \eqref{eq:critical-arithmetic-groupoid} is a countable Hausdorff,
locally compact, \'{e}tale groupoid with counting Haar system.  The formula
\[
 c_a(x,k,y)=S_ma(x)-S_na(y)
\]
is independent of the representative $(m,n)$.  The maps
\begin{equation}
\label{eq:critical-groupoid-cocycles}
 D(x,k,y)=(-k,-c_a(x,k,y)+k),
 \qquad
 \delta_{\mathcal G}(x,k,y)=-c_a(x,k,y)\log2+k\log3
\end{equation}
are additive groupoid cocycles.

For a path $p:x\to y$,
\[
 J(p)=(y,-\ell(p),x)
\]
defines an injective functor and
\[
 D(J(p))=(\ell(p),A(p)-\ell(p)),\qquad
 \delta_{\mathcal G}(J(p))=A(p)\log2-\ell(p)\log3.
\]
Therefore $[p]\mapsto1_{\{J(p)\}}$ is an injective algebra map
$\mathbb C[\mathcal E_C]\to C_c(\mathcal G_{\mathrm{arith}})$.
The two cocycles give respectively a $\mathbb T^2$ gauge action and an
$\mathbb R$ time action on both the full and reduced groupoid
$C^*$-algebras.
\end{theorem}

\begin{proof}
If $(m,n)$ is a representative and $u=C^m(x)=C^n(y)$, extending both times
by $t$ adds the same sum $S_ta(u)$ to the two terms defining $c_a$.  Any two
representatives have the same difference and admit a common extension, so
$c_a$ is well-defined.  For composable representatives $(m,n)$ and $(p,q)$,
extend them until the two middle times equal $\max\{n,p\}$.  The middle
Birkhoff sums cancel, proving additivity of both $k$ and $c_a$, hence of the
two displayed linear combinations.

Every singleton is a compact open bisection in the discrete topology, which
proves all topological and Haar assertions.  The pair $(0,\ell(p))$
represents $J(p)$ and has $c_a=-A(p)$, proving the cocycle formulas.  Source
and length determine a path of the deterministic first-return category, so
$J$ is injective.  Point masses respect the target-left product and are
linearly independent.  Multiplication by the characters of $D$ and by
$e^{it\delta_{\mathcal G}}$ gives convolution $*$-automorphisms.  Universal
norms give the full extensions and diagonal unitaries in regular
representations give the reduced extensions.
\end{proof}

\subsection{The exact 3-adic object and the missing arithmetic map}

Lines 12801--12928 introduce
\begin{equation}
\label{eq:critical-three-adic-branches}
 X=\mathbb Z_3,\qquad
 Y_a=\{y:2^ay\equiv1\pmod3\},\qquad
 g_a(x)=\frac{3x+1}{2^a},\qquad
 \beta_a(y)=\frac{2^ay-1}{3}.
\end{equation}
This is a genuine construction.

\begin{proposition}[exact branch homeomorphisms]
\label{prop:critical-three-adic-homeomorphisms}
For every $a\geq1$, $Y_a=2^{-a}+3\mathbb Z_3$ is clopen and
$g_a:X\to Y_a$ is a homeomorphism with inverse $\beta_a$.  For a word
$w=(a_1,\ldots,a_n)$,
\[
 h_w(x)=g_{a_n}\circ\cdots\circ g_{a_1}(x)
 =\frac{3^nx+B_w}{2^{A(w)}},
 \qquad
 B_w=\sum_{i=1}^n3^{n-i}2^{a_1+\cdots+a_{i-1}},
\]
and $h_w$ maps $X$ homeomorphically onto a clopen set.
\end{proposition}

\begin{proof}
The residue class is clopen because $2^a$ is a $3$-adic unit.  Both inverse
identities follow by direct substitution, and the affine word formula follows
by induction.
\end{proof}

Let $\mathcal H_3$ be the groupoid of germs of the pseudogroup generated by
these homeomorphisms and their inverses.  Each pseudogroup element is locally
affine.  Two affine maps agreeing on a nonempty $3$-adic open set have equal
slope and intercept; this proves that the germ groupoid is Hausdorff.  Its
germ bisections make it second countable, locally compact, and \'{e}tale.

\begin{theorem}[independent arithmetic-to-$3$-adic germ functor]
\label{thm:critical-arithmetic-germ-functor}
For $(x,m-n,y)\in\mathcal G_{\mathrm{arith}}$, let $p_m(x)$ and $p_n(y)$
denote the actual exponent words along the indicated positive orbits.  The
formula
\begin{equation}
\label{eq:critical-arithmetic-germ-map}
 \Xi(x,m-n,y)=
 [h_{p_m(x)}^{-1}\circ h_{p_n(y)},y]
\end{equation}
defines an injective groupoid functor
$\Xi:\mathcal G_{\mathrm{arith}}\to\mathcal H_3$.
If $u$ is the rational slope of the affine germ, then
\[
 u=3^{n-m}2^{S_ma(x)-S_na(y)},\qquad
 D=(\nu_3u,-\nu_2u-\nu_3u),\qquad
 \delta_{\mathcal G}=-\log|u|_\infty.
\]
Indeed $\nu_3u=n-m=-k$ and $\nu_2u=c_a$, so these are exactly
the coordinates $(-k,-c_a+k)$.
\end{theorem}

\begin{proof}
Adding a common future word on both sides of
\eqref{eq:critical-arithmetic-germ-map} cancels that affine map against its
inverse, so the germ is independent of $(m,n)$.  Common-middle extension
proves functoriality.  Its slope has $\nu_3(u)=n-m$ and
$\nu_2(u)=S_ma(x)-S_na(y)$.  Source, range, and $-\nu_3(u)$ recover the
triple, proving injectivity and the cocycle formulas.
\end{proof}

The raw recovery on $\ell^2(\mathcal O)$ at lines 12995--13042 is false.
Indeed $g_1(1)=2\notin\mathcal O$.  Under the printed formula for $V_1$, the
equation $\beta_1(y)=1$ has the unique solution $y=2$, so
$V_1\delta_1=0$ and $V_1^*V_1\ne1$.  The subsequent claims that the old path
algebra is exactly recovered, that an unspecified semigroup crossed product
equals $C^*(\mathcal H_3)$, and that prefixing defines the printed
endomorphism therefore do not follow.  The exact replacement is the functor
$\Xi$; because $\mathbb Z_3$ is nondiscrete, it is not a point-mass embedding
into $C_c(\mathcal H_3)$.

The prefix defect can be proved before any completion.  In the free-monoid
algebra let $L_q[p]=[qp]$.  For nonempty $q$,
\[
 L_q(1\cdot1)=[q],\qquad L_q(1)L_q(1)=[qq],
\]
and these are distinct basis vectors.  Thus left prefixing is not an algebra
endomorphism.  A genuine endomorphism would require a different declared map
and a proof that it respects every covariance relation.

\subsection{The two torus coordinates and the actual 1-motive squares}

The raw text uses the degree character $(\ell,A-\ell)$ and the arithmetic
character $(A,-\ell)$ as if they were the same coordinates.  They are related
by a specific unimodular torus isomorphism.

\begin{theorem}[independent torus and Kummer repair]
\label{thm:critical-torus-kummer-repair}
Let
\[
 T_{\mathrm{ar}}=(\mathbf G_m)^2_{x,y},\qquad
 T_{\deg}=(\mathbf G_m)^2_{\lambda,\mu},
\]
and
\[
 \Phi(x,y)=(x/y,x),\qquad
 \Phi^{-1}(\lambda,\mu)=(\mu,\mu/\lambda).
\]
For a path $p$ put
\[
 \chi_p(x,y)=x^{A(p)}y^{-\ell(p)},\qquad
 \psi_p(\lambda,\mu)=\lambda^{\ell(p)}\mu^{A(p)-\ell(p)}.
\]
Then $\psi_p\Phi=\chi_p$, with lattice matrices
\[
 \binom{A}{-\ell}
 =\begin{pmatrix}1&1\\ -1&0\end{pmatrix}
 \binom{\ell}{A-\ell},
 \qquad
 \binom{\ell}{A-\ell}
 =\begin{pmatrix}0&-1\\ 1&1\end{pmatrix}
 \binom{A}{-\ell}.
\]

Over $\mathbb Q$, define
\[
 M_{\mathrm{ar}}=[\mathbb Z\to T_{\mathrm{ar}}],
 \quad n\mapsto(2^n,3^n),
 \qquad
 K(\alpha_p)=[\mathbb Z\to\mathbf G_m],
 \quad n\mapsto\alpha_p^n,
\]
where $\alpha_p=2^{A(p)}3^{-\ell(p)}$.  Then
$(\mathrm{id},\chi_p):M_{\mathrm{ar}}\to K(\alpha_p)$ is a morphism of
$1$-motives and, for $\gamma(s)=(2^s,3^s)$,
\[
 \int_\gamma\chi_p^*\frac{dz}{z}
 =A(p)\log2-\ell(p)\log3.
\]
The path fixes the logarithm branch.  The same square with
$\chi_{p,k}(x,y)=x^{-A_k}y^{k-1}$ gives the period
$\log(3^{k-1}2^{-A_k})$.
\end{theorem}

\begin{proof}
Substitution proves the torus inverse and character identity; multiplication
of the integral matrices gives the identity.  The square of group schemes
commutes because
\[
 \chi_p(2^n,3^n)=(2^{A(p)}3^{-\ell(p)})^n.
\]
Finally
$\chi_p^*(dz/z)=A(p)\,dx/x-\ell(p)\,dy/y$, and integration on the declared
path proves the period formula.  The prefix calculation is identical.
\end{proof}

The Cartier dual specialization is
\[
 M_{\mathrm{ar}}^\vee=[\mathbb Z^2\to\mathbf G_m],\qquad
 (r,s)\mapsto2^r3^s.
\]
Thus the alternate raw object
$[\mathbb Z^2\to\mathbf G_m]$, $(m,n)\mapsto2^m3^{-n}$, is the same dual
after the declared basis involution $(m,n)\mapsto(m,-n)$.  It is not a
replacement for the commutative square just proved.

\subsection{The arctangent pair and the Nori edge direction}

Over $K=\mathbb Q(i)$, put
$U_K=\mathbf P^1_K\setminus\{i,-i\}$.  The Cayley map and inverse are
\[
 \kappa(z)=\frac{1+iz}{1-iz},
 \qquad
 \kappa^{-1}(w)=\frac{w-1}{i(w+1)}.
\]

\begin{proposition}[independent Cayley pair repair]
\label{prop:critical-cayley-nori-repair}
The map $\kappa$ is an isomorphism $U_K\simeq\mathbf G_{m,K}$.  For
$r\in\mathbb Q_{>0}$ and $\alpha(r)=(1+ir)/(1-ir)$ it is an isomorphism of
pairs
\[
 (U_K,\{0,r\})\overset\sim\longrightarrow
 (\mathbf G_{m,K},\{1,\alpha(r)\}),
 \qquad
 \kappa^*\frac{dw}{w}=\frac{2i\,dz}{1+z^2}.
\]
The resulting Nori functoriality edge is contravariant:
\[
 (\mathbf G_{m,K},\{1,\alpha(r)\},1)\longrightarrow
 (U_K,\{0,r\},1).
\]
Over $\mathbb Q$, the same integral is a $1$-period of
$(\mathbf P^1\setminus V(z^2+1),\{0,r\},1)$.
\end{proposition}

\begin{proof}
Solving the fractional-linear equation proves the inverse.  The deleted
points map to $0$ and $\infty$, while $0\mapsto1$ and
$r\mapsto\alpha(r)$.  Logarithmic differentiation proves the differential
identity.  The Nori edge direction is exactly the contravariant
functoriality direction in Definition~9.1.1 of
\cite{HuberMullerStach2025}.  The rational form $dz/(1+z^2)$ is regular away
from $V(z^2+1)$.  At infinity the coordinate $t=1/z$ gives
\[
 \frac{dz}{1+z^2}=-\frac{dt}{1+t^2},
\]
so it is regular there as well.  Hence the rational relative curve pair
supplies the period.
\end{proof}

The raw curve at lines 11515--11518 deletes infinity unnecessarily.  Under
$\kappa$, deleting infinity also deletes $-1$ from $\mathbf G_m$, so it is
not the full pair isomorphism printed above.

The raw assertion at lines 11574--11588 of a path-indexed Nori diagram is
not constructed.  It omits the degree $i$, switches the attached Nori
volume's convention $(k,\mathbb Q)$-Vect to $(\mathbb Q,k)$-Vect, and gives
neither maps of pairs nor nested triples.  Concatenating singular paths gives
a relation among homology classes, not a map of ambient algebraic pairs.

There is nevertheless an exact multiplication map.  For composable $p,q$,
\[
 \chi_{q\circ p}=\chi_q\chi_p,\qquad
 \alpha_{q\circ p}=\alpha_q\alpha_p.
\]
Put
\[
 N_{p,q}=[\mathbb Z\to\mathbf G_m^2],\qquad
 n\mapsto(\alpha_p^n,\alpha_q^n).
\]
Then
\[
 M_{\mathrm{ar}}
 \xrightarrow{(\mathrm{id},(\chi_p,\chi_q))}N_{p,q}
 \xrightarrow{(\mathrm{id},m)}K(\alpha_p\alpha_q),\qquad
 m(z_1,z_2)=z_1z_2,
\]
is a pair of commutative $1$-motive squares.  The map of pairs
\[
 m:(\mathbf G_m^2,\{(1,1),(\alpha_p,\alpha_q)\})
 \to(\mathbf G_m,\{1,\alpha_p\alpha_q\})
\]
gives the genuine contravariant Nori edge
\[
 (\mathbf G_m,\{1,\alpha_p\alpha_q\},1)
 \xrightarrow{m^*}
 (\mathbf G_m^2,\{(1,1),(\alpha_p,\alpha_q)\},1),
\]
and
\[
 m^*\frac{dz}{z}=\operatorname{pr}_1^*\frac{dz}{z}
 +\operatorname{pr}_2^*\frac{dz}{z}
\]
proves period additivity.  This is the exact replacement for the missing
``concatenation arrow.''

Finally, for $r>0$,
\[
 \operatorname{Im}\frac{1+ir}{1-ir}=\frac{2r}{1+r^2}>0,
\]
whereas $2^{A(p)}3^{-\ell(p)}$ is positive rational.  The two marked Kummer
points are therefore unequal; no direct edge between those literal pairs is
obtained by identifying them.

\subsection{Status of the remaining assertions}

The status of the raw interval is now exact:
\begin{itemize}
\item the bidegree, torus action, finite residue maps, full-shift pressure,
  logarithmic formulas, and arctangent integrals survive at the proved scopes
  stated here and in the preceding sections;
\item the discrete arithmetic groupoid, the $3$-adic germ groupoid, the
  faithful functor $\Xi$, the torus coordinate isomorphism, the Kummer
  squares, and the Cayley and multiplication maps are separately authored
  repairs proved in this edition;
\item the printed $\ell^2(\mathcal O)$ isometry and exact path-algebra
  recovery are disproved by $g_1(1)=2$;
\item the unnamed crossed-product equality, prefix endomorphism, equilibrium
  classification, cyclic-homology chain, motivic coaction, absolute-arithmetic
  interpretation, and local-factor package are not constructed by the raw
  formulas and are not admitted here.
\end{itemize}

\begin{nonclaim}
The exact maps above do not prove amenability, equality of full and reduced
groupoid completions, existence or classification of equilibrium states, a
universal path-indexed Nori category, a motivic Galois action on the path
algebra, or a Collatz endpoint theorem.  These are precise nonclaims about
this audited interval, not assertions that no such mathematics can be
constructed elsewhere.
\end{nonclaim}
